\documentclass[12pt]{article}
\usepackage[utf8]{inputenc}
\usepackage{graphicx}
\usepackage{amsmath}
\usepackage{amssymb}
\usepackage{mathrsfs}
\usepackage{pgfornament}
\usepackage[many]{tcolorbox}
\usepackage[margin = 0.5in]{geometry}
\usepackage{collectbox}
\usepackage{mdframed}
\usepackage{xcolor}
\usepackage{mathtools}
\DeclarePairedDelimiter\ceil{\lceil}{\rceil}
\DeclarePairedDelimiter\floor{\lfloor}{\rfloor}
\newcommand\norm[1]{\left\lVert#1\right\rVert}
\newcommand\textbox[1]{%
  \parbox{.333\textwidth}{#1}%
}

\linespread{2}

\makeatletter
\newcommand{\mybox}{%
    \collectbox{%
        \setlength{\fboxsep}{2pt}%
        \fbox{\BOXCONTENT}%
    }%
}
\makeatother

\newcommand{\innerprod}[2]{\left<#1, #2\right>}

\usepackage{listings}
\usepackage{color}

\definecolor{dkgreen}{rgb}{0,0.6,0}
\definecolor{gray}{rgb}{0.5,0.5,0.5}
\definecolor{mauve}{rgb}{0.58,0,0.82}

\lstset{frame=tb,
  language=R,
  aboveskip=3mm,
  belowskip=3mm,
  showstringspaces=false,
  columns=flexible,
  basicstyle={\small\ttfamily},
  numbers=none,
  numberstyle=\tiny\color{gray},
  keywordstyle=\color{blue},
  commentstyle=\color{dkgreen},
  stringstyle=\color{mauve},
  breaklines=true,
  breakatwhitespace=true,
  tabsize=3
}
\title{MAT 523 HW 8}
\begin{document}

\hfill April 03, 2019
\begin{center}
\begin{large}
\textbf{MAT 524: Functions of a Real Variable II \\
\vspace{-5mm}
Homework VIII \\}
\end{large}
\vspace{-2mm}
Nolan R. H. Gagnon \\
\end{center}
\vspace{-0.5cm}
\noindent \textbf{{\huge [}{\Large [}} Problem One \textbf{{\Large ]}{\huge ]}}{\Large \textbf{:}}

\vspace{3mm}

\noindent\mybox{\colorbox{yellow}{\textbf{Problem Statement}}}\hspace{3mm}\hrulefill

\noindent Let $f(x) = \frac{1}{2} - x$ on $[0,1)$, viewed as an element of $L^2([0,1])$. Show that $\hat{f}(0)=0$ and $\hat{f}(n) = \frac{1}{2\pi i n}$ if $n\neq 0$. Apply Parseval's identity to compute $\zeta(2)$: 
$$\sum\limits_{n = 1}^{\infty} \frac{1}{n^2}= \frac{\pi^2}{6}.$$

\noindent\mybox{\colorbox{yellow}{\textbf{Solution}}}\hspace{3mm}\hrulefill

\noindent Note that the set $\mathcal{B} = \left\lbrace e^{2\pi i n x} \hspace{2mm}|\hspace{2mm} n \in \mathbb{Z}\right\rbrace$ is an orthonormal basis for $L^2([0,1]).$ Using this set, we have
\begin{align}
    \hat{f}(0) &= \innerprod{f(x)}{e^0} \\
    &= \int_0^1 \left(\frac{1}{2} - x\right) dx \\
    &= \left[\frac{1}{2}x - \frac{x^2}{2}\right]_0^1 \\
    &= \frac{1}{2} - \frac{1}{2} \\
    &= 0.
\end{align}
For $n\neq 0$, we have
\begin{align}
    \hat{f}(n) &= \innerprod{f(x)}{e^{2\pi i n x}} \\
    &= \int_0^1 \left(\frac{1}{2} - x\right)\overline{e^{2\pi i n x}} dx \\
    &= \int_0^1 \left(\frac{1}{2} - x\right)e^{-2\pi i n x} dx \\
    &= \int_0^1 \frac{1}{2}e^{-2\pi i n x} dx - \int_0^1 xe^{-2\pi i n x}dx.
\end{align}
For the first integral in Equation (9), we have
\begin{align}
    \int_0^1 \frac{1}{2}e^{-2\pi i n x} dx &= \frac{1}{-4\pi i n}\left[ e^{-2\pi i n} - 1\right] \\
    &= 0.
\end{align}
To compute the second integral from Equation (9), we use integration by parts:
\begin{align}
    -\int_0^1 x e^{-2\pi i n x} dx &= -\left[\frac{-x}{2\pi i n}e^{-2\pi i n x}\right]_0^1 - \int_0^1 \frac{1}{2\pi i n}e^{-2\pi i n x} dx \\
    &= \frac{1}{2\pi i n} - \left[\frac{1}{-4\pi^2 i^2 n^2}e^{-2\pi i n x}\right]_0^1 \\
    &= \frac{1}{2\pi i n}. 
\end{align}
Thus, $\hat{f}(n) = \frac{1}{2\pi i n}.$ 

Now, since $\mathcal{B}$ is an orthonormal basis for $L^2([0,1])$, Bessel's inequality becomes Parseval's identity. So
\begin{align}
    \sum\limits_{n=-\infty}^{\infty} \left|\hat{f}(n)\right|^2 &= 2\sum\limits_{n=1}^{\infty} \left|\frac{1}{2\pi i n}\right|^2 \\
    &= \sum\limits_{n=1}^{\infty}\frac{1}{2\pi^2 n^2} \\
    &= \norm{f(x)}^2 \\
    &= \int_0^1 \left|\frac{1}{2}-x\right|^2 dx \\
    &= \int_0^{0.5} \left(0.5 -x\right)^2 dx + \int_{0.5}^{1} \left(x - 0.5\right)^2 dx \\
    &= \left[0.25x -\frac{x^2}{2}+\frac{x^3}{3}\right]_0^{0.5} + \left[\frac{x^3}{3} - \frac{x^2}{2} + 0.25x\right]_{0.5}^{1} \\
    &= \frac{1}{12}.
\end{align}
Therefore, $$\sum\limits_{n=1}^{\infty} \frac{1}{n^2} &= \frac{\pi^2}{6}.$$

\hfill$\blacksquare$

\vfill

\begin{center}
\pgfornament[width = 50mm, color = black, symmetry = h]{83}\\
\vspace{2mm}
\vspace{-0.65cm}
\hrulefill \\
\vspace{-0.95cm}
\hrulefill
\end{center}

\end{document}
